% ID: SEL-D0060-OA-E02
% ESTADO: validada. S060.
\textbf{a)} En $6$ ambos tramos coinciden en $7/2$. Sus derivadas son $t/4-1$ y $1/2$. Las dos pendientes laterales en $6$ valen $1/2$, por lo que $B'(6)=1/2$.

\textbf{b)} La derivada del primer tramo cambia de negativa a positiva en $4$ y el segundo tramo sigue creciendo. El mínimo es $B(4)=3$, es decir, $3000$ euros al cuarto mes.

\textbf{c)} La parábola pasa por $(0,5)$, $(4,3)$ y $(6,7/2)$. Desde este último punto parte el segmento hasta $(12,13/2)$. Comparando los dos extremos del dominio, el máximo es $6500$ euros al cabo de $12$ meses.
\begin{center}
\setlength{\unitlength}{0.9mm}
\begin{picture}(96,57)(-10,-6)
\qbezier(0,0)(40,0)(80,0)
\qbezier(0,0)(0,24)(0,48)
\qbezier(0,35)(18,14)(36,24.5)
\qbezier(36,24.5)(54,35)(72,45.5)
\put(81,0){$t$}
\put(2,48){$B(t)$}
\put(2,36){\scriptsize $(0,5)$}
\put(19,15){\scriptsize $(4,3)$}
\put(36,18){\scriptsize $(6,7/2)$}
\put(72,46){\scriptsize $(12,13/2)$}
\end{picture}
\end{center}
